% SPDX-License-Identifier: GPL-3.0-or-later
% Copyright (C) 2026 Alessandro Greco
\chapter{Защита от DoS и управление ресурсами}\label{cap:15-dos}
\aggiornatoal{23 сентября 2026 г.}{коммит torspec \texttt{8e1f245}; tor 0.4.9.13; dos-spec, hspow-spec}

\begin{obiettivi}
К концу этой главы вы сможете:
\begin{itemize}[nosep]
  \item объяснять, почему атака типа «отказ в обслуживании» против Tor одновременно угрожает анонимности;
  \item описывать, как tor защищает память реле, ресурсы процессора и пропускную способность;
  \item понимать ограничения числа соединений и цепей по IP-адресам и роль параметров консенсуса;
  \item объяснять, как onion-сервисы защищаются от массовых запросов введения;
  \item описывать механизм доказательства выполнения работы (PoW) версии 1 (Equi-X и Blake2b) и управление вычислительной сложностью;
  \item вызывать и наблюдать защитные механизмы реле и onion-сервиса в лабораторной сети (лабораторная работа~9).
\end{itemize}
\end{obiettivi}

Открытая анонимная сеть, поддерживаемая добровольцами, естественным образом становится целью тех, кто хочет сделать её непригодной к использованию. В этой главе рассматриваются механизмы защиты Tor от \emph{отказа в обслуживании} (denial of service, DoS): попыток исчерпать ресурсы реле или onion-сервиса, чтобы ими не могли пользоваться другие. Конкретные инциденты последних лет, включая волну атак на onion-сервисы в 2022--2023 годах, рассматриваются в \cref{cap:27-casi-rete}. Здесь мы сосредоточимся на механизмах.

% ---------------------------------------------------------------------------
\section{Почему DoS --- проблема анонимности}\label{sec:dos-anonimato}

В обычной сети DoS --- \enquote{лишь} проблема доступности. В Tor это также оружие против анонимности~\spec{dos-spec}{Overview}:
\begin{itemize}
  \item Противник, способный по требованию вывести реле из строя, может заставить клиентов, использующих его как входной узел (guard), выбрать другой, рассчитывая, что новым входным узлом станет реле противника. Jansen, Tschorsch, Johnson и Scheuermann показали, как противник с небольшими ресурсами может исчерпать память реле и использовать эту возможность для деанонимизации onion-сервисов: \emph{атака Sniper Attack}~\cite{jansen2014sniper}.
  \item Замедление или недоступность честных реле перенаправляет трафик к враждебным реле, увеличивая долю наблюдаемого ими трафика.
  \item Атака, выборочно вызывающая отказы цепей, может служить сигналом, например для определения цепей, проходящих через конкретное реле.
\end{itemize}
Поэтому защита должна не просто \enquote{поддерживать работу реле}, но и не создавать наблюдаемого поведения, которое способен использовать противник.

% ---------------------------------------------------------------------------
\section{Физические ресурсы}\label{sec:dos-risorse}

\subsection{Память}

Память --- наиболее чувствительный ресурс и защищена наиболее зрелыми механизмами~\spec{dos-spec}{Memory exhaustion}. В Linux процесс, запросивший слишком много памяти, обычно не получает ошибку сразу. Система обещает выделить память и может позднее завершить процесс, когда тот действительно попытается её использовать. Поэтому ждать ошибки выделения памяти бессмысленно. tor устанавливает собственный лимит, \opt{MaxMemInQueues}, задаваемый оператором или вычисляемый как доля системной памяти. В него учитывается всё, что потребляет память: ячейки в очередях цепей, буферы соединений, состояние сжатия, кэшированные дескрипторы onion-сервисов и DNS-запросы.

При превышении лимита tor освобождает память в два этапа. Сначала из кэшей, потребляющих более 20\,\% общего объёма, удаляются самые старые записи, пока каждый такой кэш не опустится ниже 10\,\%. Если использование всё ещё выше 90\,\% лимита, tor закрывает цепи и соединения в порядке возраста наиболее старых поставленных в очередь данных, начиная с тех, где данные стоят в очереди дольше всего. Цель --- затронуть участников, которые накапливают данные, не читая их, что характерно для подобных атак.

\subsection{Процессор}

Нагрузка на процессор --- одно из слабых мест реализации на C, обрабатывающей сетевые события в одном главном потоке. Наиболее дорогие операции, криптографические рукопожатия при создании цепей (\cref{cap:09-circuiti}), делегируются пулу рабочих потоков и ставятся в очередь. Но при слишком большом числе рукопожатий очередь растёт, и легитимные цепи завершаются по тайм-ауту. Tor не измеряет загрузку CPU напрямую. Он защищает CPU косвенно, ограничивая операции, которые его расходуют, включая создание цепей и обработку запросов введения для onion-сервисов.

\subsection{Пропускная способность}

Операторы могут ограничивать пропускную способность с помощью параметров вроде \opt{RelayBandwidthRate}, но такие лимиты прежде всего не позволяют выйти за условия канала связи и сами по себе не являются защитой. От аномального расхода пропускной способности Tor защищается механизмами управления потоком и перегрузкой (\cref{cap:11-stream}) и справедливым планированием ячеек между цепями.

% ---------------------------------------------------------------------------
\section{Ограничения по IP-адресам}\label{sec:dos-limiti}

Самая эффективная защита от простых атак с массовой отправкой запросов проста. Входной узел (guard) видит IP-адрес подключающегося клиента и может ограничить объём запросов от каждого адреса. Это единственное место сети, где применим такой лимит, поскольку последующие переходы намеренно не знают источник трафика. Tor применяет к клиентам, но не к аутентифицированным реле, три семейства ограничений, настраиваемых опциями \texttt{DoS\ldots}:
\begin{description}
  \item[Соединения] (\texttt{DoSConnection\ldots}) Максимальное число одновременных TLS-соединений и максимальная скорость новых соединений с одного адреса. Превысившие пороги соединения закрываются.
  \item[Цепи] (\texttt{DoSCircuitCreation\ldots}) \emph{корзина токенов (token bucket)} для каждого адреса: каждая новая цепь расходует один токен. Токены пополняются с фиксированной скоростью до предельного запаса (\emph{burst}). Адрес, исчерпавший токены \emph{и} удерживающий как минимум заданное число соединений, помечается; последующие запросы цепей отклоняются на определённый период.
  \item[Одношаговые запросы к пункту встречи] (\texttt{DoSRefuseSingleHopClientRendezvous}) Отклоняет запросы к пунктам встречи по цепям из одного перехода, например от шлюзов Tor2web, позволяющих обращаться к onion-сервисам без анонимности клиента. Такие запросы представляют дешёвый способ генерировать нагрузку.
\end{description}

\subsection{Код и консенсус}

В коде tor эти механизмы \textbf{отключены} по умолчанию (\file{DOS\_CC\_ENABLED\_DEFAULT 0} в \file{src/core/or/dos.h}). Общесетевые настройки поступают из параметров консенсуса. Публичный консенсус от 23 сентября 2026 года содержит:
\begin{center}\small
\begin{tabular}{lrr}
\toprule
Параметр & Консенсус & Значение в коде по умолчанию \\
\midrule
\texttt{DoSCircuitCreationEnabled} & 1 & 0 \\
\texttt{DoSCircuitCreationMinConnections} & 2 & 3 \\
\texttt{DoSCircuitCreationRate} & 2 & 3 \\
\texttt{DoSCircuitCreationBurst} & 60 & 90 \\
\texttt{DoSConnectionEnabled} & 1 & 0 \\
\texttt{DoSConnectionMaxConcurrentCount} & 50 & 100 \\
\texttt{DoSRefuseSingleHopClientRendezvous} & 1 & --- \\
\bottomrule
\end{tabular}
\end{center}
Сетевые значения строже значений по умолчанию в коде: они отражают годы атак и последующих корректировок. Возможность менять их через консенсус за считанные часы без обновления тысяч реле --- один из наиболее ценных инструментов Tor Project во время атаки.

\begin{attenzione}[Недостатки ограничений по IP-адресам]
Противник со множеством IP-адресов, например ботнет, может обходить ограничения по IP-адресам. С другой стороны, слишком строгие лимиты затрагивают легитимных пользователей, делящих один адрес, например за NAT мобильного оператора или университета. Значения всегда являются компромиссом.
\end{attenzione}

% ---------------------------------------------------------------------------
\section{Onion-сервисы под атакой}\label{sec:dos-onion}

Onion-сервисы особенно уязвимы. Каждый запрос введения обходится клиенту дёшево: один \texttt{INTRODUCE1} (\cref{sec:os-rendezvous}). Для сервиса тот же запрос дорог: необходимо расшифровать сообщение, выполнить рукопожатие и построить многоузловую цепь к указанному пункту встречи. Возникает \textbf{усиление нагрузки}: имея сравнительно небольшие ресурсы, противник заставляет сервис и все реле его цепей выполнять значительно больше работы. Сервис не может блокировать IP-адреса атакующих, поскольку никогда их не видит.

\subsection{Ограничения в пунктах введения}

Первая защита --- ограничение частоты запросов в пунктах введения. С помощью \texttt{HiddenService\-EnableIntroDoS\-Defense} и связанных \texttt{\ldots RatePerSec} и \texttt{\ldots BurstPerSec} сервис через расширение \file{DOS\_PARAMS} в \file{ESTABLISH\_INTRO} сообщает каждому пункту введения, сколько сообщений \texttt{INTRODUCE2} в секунду готов принимать. Это защищает сервис, но не отличает атакующего от легитимного пользователя. При атаке могут быть отклонены запросы и атакующего, и законного пользователя.

\subsection{Доказательство выполнения работы}

Более новый механизм, предложенный в \proposta[Предложении]{327} и доступный с tor 0.4.8.4 (август 2023 года), --- \textbf{доказательство выполнения работы} (proof of work, PoW)~\spec{hspow-spec}{Motivation}. Клиенты должны доказать, что выполнили определённый объём вычислений перед обслуживанием; запросы, подтверждающие большее вычислительное усилие, получают приоритет. Несколько секунд вычислений приемлемы для легитимного пользователя. Для атакующего, посылающего тысячи запросов в секунду, умножение этой стоимости на каждый запрос может стать запретительным.

Механизм состоит из четырёх частей~\spec{hspow-spec}{Common protocol}:
\begin{enumerate}
  \item \textbf{Параметры дескриптора.} Во втором слое дескриптора (\cref{sec:os-descrittore}) сервис публикует строку \texttt{pow-params} со схемой (\texttt{v1}), случайным 32-байтовым \emph{семенем}, \emph{рекомендуемым усилием} и сроком действия семени. Семя обновляется каждые 105--120 минут.
  \item \textbf{Клиент решает задачу.} Выбрав усилие $E$, клиент ищет решение Equi-X для задачи
  \[ \texttt{"Tor hs intro v1\textbackslash 0"} \,\|\, A' \,\|\, \mathit{seed} \,\|\, \mathit{nonce} \,\|\, E, \]
  где $A'$ --- ослеплённый ключ сервиса. Одного решения Equi-X недостаточно. Клиент также вычисляет 32-битное значение хеша Blake2b $R$ от задачи и решения; решение принимается только если $R \cdot E \le 2^{32} - 1$. Иначе меняется одноразовое число (nonce) и попытка повторяется.
  \item \textbf{Доказательство в сообщении введения.} Клиент включает одноразовое число nonce, усилие, первые 4 байта семени и решение в расширение зашифрованной части \texttt{INTRODUCE1}, доступной только сервису.
  \item \textbf{Очередь приоритетов.} Сервис проверяет доказательство --- значительно быстрее, чем его находят, --- и помещает запрос в очередь, упорядоченную по усилию. Сначала обслуживаются запросы с наибольшим усилием; запросы без доказательства имеют нулевое усилие. При заполнении очереди сервис сразу отбрасывает половину с наименьшим усилием.
\end{enumerate}

\begin{figure}[htbp]
\centering
% SPDX-License-Identifier: GPL-3.0-or-later
% Copyright (C) 2026 Alessandro Greco
%% Русское альтернативное описание EPUB.
\altfig{Дополнительная защита PoW онион-сервиса. Если она включена, сервис публикует pow-params с семенем и рекомендуемым усилием; клиент выбирает фактическое E, ищет nonce и решение Equi-X, также удовлетворяющее R умножить на E не больше 2^32−1. Доказательство передаётся в зашифрованном расширении INTRODUCE1 через точку введения как INTRODUCE2. Сервис проверяет семя, повтор nonce, правильность Equi-X и усилие; неверная проверка приводит к отклонению. Проверенные запросы становятся в очередь по убыванию усилия, запросы без доказательства имеют E=0. При переполнении реализация отбрасывает половину с меньшим усилием. Числа 900, 400, 120, 40 и 0 приведены как пример.}
% Доказательство работы онион-сервиса: от дескриптора к очереди приоритетов
\begin{tikzpicture}[font=\sffamily\scriptsize,
  box/.style={draw,rounded corners=2pt,align=center,minimum height=1cm}]
  \node[box,fill=torgrigio!8,text width=2.2cm] (d) at (0,0) {дескриптор:\\\texttt{pow-params}\\семя, рекомендованное усилие};
  \node[box,fill=torviola!10,text width=3.2cm] (c) at (3.8,0) {клиент: выбирает $E$, ищет nonce и\\решение Equi-X с\\$R \cdot E \le 2^{32}-1$};
  \node[box,fill=torgrigio!8,text width=2.2cm] (i) at (8.4,0) {точка\\введения\\(ограничение частоты)};
  \node[box,fill=torverde!12,text width=2.2cm] (s) at (8.4,-2.2) {сервис:\\проверяет доказательство\\(быстро)};
  \draw[->] (d) -- (c);
  \draw[->] (c) -- node[below,pos=0.45] {\texttt{INTRODUCE1}} (i);
  \draw[->] (i) -- node[right] {\texttt{INTRODUCE2}} (s);
  \foreach \k/\e/\c in {0/900/torverde!45,1/400/torverde!35,2/120/torverde!25,3/40/torverde!15,4/0/torgrigio!10}
    \node[draw,minimum width=1.1cm,minimum height=0.55cm,fill=\c] (q\k) at (1.2+1.1*\k,-2.2) {$E{=}\e$};
  \draw[->] (s) -- (q4.east);
  \node[above=0.1cm of q2] {очередь приоритетов, по убыванию усилия (без доказательства: $E=0$)};
  \node[below=0.1cm of q0,align=center] {обслуживается первым};
  \node[below=0.1cm of q4,align=center] {очередь заполнена: отбрасывается\\половина с наименьшим усилием};
\end{tikzpicture}

\caption{Доказательство выполнения работы для onion-сервиса. Клиент читает семя и рекомендуемое усилие из дескриптора, ищет решение, проходящее проверку, и отправляет его в зашифрованной части \texttt{INTRODUCE1}. Сервис проверяет решение и упорядочивает запросы по усилию.}
\label{fig:pow}
\end{figure}

\begin{approfondimento}[Почему Equi-X?]
Equi-X, разработанный tevador для этой задачи, --- \emph{асимметричная} задача: поиск решения требует памяти и значительных вычислений, а проверка очень дешева. Он основан на HashX --- семействе хеш-функций, реализации которых меняются с каждым семенем и задействуют возможности параллельной работы современных процессоров. Это усложняет ускорение с помощью специализированного оборудования и ограничивает преимущество атакующего над пользователем обычного компьютера. Существуют две реализации: C-версия в tor и Rust-версия для Arti, совместимость которых поддерживают автоматические тесты. C-код распространяется под LGPL, поэтому для включения tor должен собираться с \texttt{--enable-gpl}. Сборка без этой опции не поддерживает PoW: \cmd{tor --list-modules} показывает \texttt{pow: no}, а не \texttt{pow: yes}, который мы увидим в \cref{cap:17-torrc}. Это не типичный случай: и лабораторный образ, и пакет Debian \texttt{tor} собраны с этой опцией и содержат модуль.
\end{approfondimento}

\subsection{Линейная шкала усилия}

Зачем нужна проверка Blake2b? Стоимость поиска отдельного решения Equi-X приблизительно постоянна. Случайное решение удовлетворяет $R \cdot E \le 2^{32}-1$ с вероятностью примерно $1/E$, поэтому клиент в среднем должен попробовать $E$ решений. Стоимость растёт \emph{линейно} с $E$, что позволяет плавно её регулировать. Многие PoW-схемы вместо этого требуют определённое число нулевых старших битов хеша, удваивая стоимость с каждым дополнительным битом: такая шкала слишком груба для управления с обратной связью. Кроме того, усилие объявляется \emph{до} поиска решения и входит в задачу. Клиент не может сначала случайно получить особенно удачный результат, а затем заявить большее усилие, чем реально затрачен.

\subsection{Управление усилием}

Рекомендуемое усилие не является постоянным. Сервис начинает с нуля, то есть защита неактивна. Каждые пять минут он обновляет рекомендацию по состоянию очереди: увеличивает значение, если очередь оставалась полной или пришлось отбрасывать запросы с рекомендуемым уровнем усилия, и уменьшает иначе. Логика управления вдохновлена управлением перегрузкой TCP: аддитивное увеличение и мультипликативное уменьшение. Клиенты, со своей стороны, повышают усилие после истечения времени ожидания подключения. Получается система обратной связи из двух контуров, где более настойчивый пользователь может предложить большее усилие, чтобы продвинуться в очереди.

% ---------------------------------------------------------------------------
\section{Лабораторная работа: вызов защитных механизмов}\label{sec:lab-dos}

\begin{laboratorio}[9]{Ограничения против DoS и доказательство выполнения работы в лабораторной сети}
\textbf{Цель.} Включить на реле ограничения создания цепей и превысить их; включить доказательство выполнения работы для onion-сервиса и прочитать его параметры; изучить проверку уровня усилия через Blake2b с помощью \file{labs/ru/код/torlab/pow.py}.\\
\textbf{Предварительные условия.} Лабораторные работы~1--8; работающая лабораторная сеть с образом, собранным с \texttt{--enable-gpl}.\\
\textbf{Оценочное время.} 2 часа.
\end{laboratorio}

\begin{attenzione}[Только в собственной сети]
Эта лабораторная работа намеренно создаёт аномальный трафик. Выполняйте её исключительно в частной сети книги. Повторение тех же действий против реальных реле или сервисов является атакой, причиняющей вред другим людям, и во многих юрисдикциях незаконно (\cref{cap:25-legale}).
\end{attenzione}

\paragraph{Шаг 1: включить ограничения на реле.} Добавляем очень низкие лимиты в \file{torrc} реле \texttt{test003r}, чтобы они быстро сработали, и отправляем сигнал \texttt{HUP} для перезагрузки конфигурации:
\begin{lstlisting}[language=torrc,numbers=none]
DoSCircuitCreationEnabled 1
DoSCircuitCreationMinConnections 2
DoSCircuitCreationRate 1
DoSCircuitCreationBurst 10
DoSCircuitCreationDefenseType 2
DoSCircuitCreationDefenseTimePeriod 60 seconds
\end{lstlisting}
Тип защиты 2 означает \enquote{отклонять новые цепи}.

\paragraph{Шаг 2: превысить лимиты.} Открываем один или два канала к реле и быстро запрашиваем 30 цепей через \cella{CREATE\_FAST}. Чтобы каждый запуск имел отдельные счётчики, соединения исходят с разных адресов обратной петли (127.0.0.2, 127.0.0.3), используя новый параметр \texttt{source\_address} функции \file{open\_channel}:
\begin{lstlisting}[numbers=none]
источник 127.0.0.2, 2 канала: 9 успешно, 21 DESTROY(5); первый отказ: 10
источник 127.0.0.3, 1 канал: 30 успешно
источник 127.0.0.1, 2 канала: 9 успешно, 21 DESTROY(5); первый отказ: 10
\end{lstlisting}
Следуют три наблюдения:
\begin{itemize}
  \item Реле отклоняет запросы с помощью ячейки \cella{DESTROY} с кодом причины 5, \texttt{RESOURCELIMIT}.
  \item При \emph{burst} 10 успешно создаются девять цепей, а не десять. Код (\file{cc\_has\_exhausted\_circuits} в \file{src/core/or/dos.c}) помечает клиента сразу, когда запас токенов исчерпан, поэтому десятый запрос уже отклоняется.
  \item При одном канале ограничение не срабатывает: защита требует не менее \texttt{MinConnections} одновременных соединений. Это не наказывает обычного клиента, который обычно держит только одно соединение со своим входным узлом (guard).
\end{itemize}
Мы предполагали, что соединения с 127.0.0.1, адреса лабораторных реле, будут исключением. Эксперимент опроверг предположение, а \file{dos.c} подтверждает отсутствие такого исключения для цепей. В конце восстанавливаем исходный \file{torrc}.

\begin{sloppypar}
\paragraph{Шаг 3: доказательство выполнения работы onion-сервиса.} После строк \texttt{HiddenService\ldots} в \file{torrc} узла \texttt{test009h} добавляем:
\end{sloppypar}
\begin{lstlisting}[language=torrc,numbers=none]
HiddenServicePoWDefensesEnabled 1
\end{lstlisting}
После сигнала \texttt{HUP} сервис публикует новый дескриптор. Его загрузка и расшифрование с помощью лабораторной работы~8 показывает во втором слое:
\begin{lstlisting}[numbers=none]
pow-params v1 tx2Mtmzw5GHgWjcDoivTvWzwQPNZONwMwSMddOc+ZVQ= 0 2026-09-23T18:20:53
\end{lstlisting}
Строка содержит схему \texttt{v1}, семя в кодировке base64, рекомендуемое усилие \textbf{ноль} (очередь пока не создавала нагрузки) и срок действия семени. Срок истекает через час и пятьдесят минут после публикации, в предусмотренном диапазоне 105--120 минут.

\paragraph{Шаг 4: линейный рост усилия.} Модуль \file{pow.py} реализует проверку Blake2b и имитирует решения Equi-X случайными байтами, имеющими для этой проверки то же распределение:
\begin{lstlisting}[language=Python]
def check_value(service_id, seed, nonce, effort, solution):
    challenge = P_STRING + service_id + seed + nonce + struct.pack(">I", effort)
    h = hashlib.blake2b(challenge + solution, digest_size=4).digest()
    return int.from_bytes(h, "big")

def valid_effort(r, effort):
    return r * effort <= UINT32_MAX
\end{lstlisting}
Среднее число попыток до успеха в 400 испытаниях для каждого значения усилия:
\begin{lstlisting}[numbers=none]
усилие     1: 1.0
усилие    10: 9.7
усилие   100: 102.3
усилие  1000: 1068.8
\end{lstlisting}
Число попыток растёт пропорционально усилию, как и предсказано. Ожидаемое значение равно усилию, поскольку каждая попытка успешна с вероятностью $1/\mathit{effort}$, а число попыток подчиняется геометрическому распределению. Стандартное отклонение сопоставимо со средним значением. Поэтому среднее по 400 испытаниям имеет типичное относительное отклонение $1/\sqrt{400} = 5\,\%$. Два запуска могут различаться на десять процентов без каких-либо изменений; при повторении мы действительно получили 969,9 вместо 1068,8 в листинге. Это полезное напоминание: выборочное среднее следует рассматривать с учётом статистической неопределённости. Ещё одна деталь: \texttt{blake2b} с \file{digest\_size=4} не равен первым четырём байтам 64-байтового результата Blake2b, поскольку длина выхода влияет на начальное состояние функции. Спецификация отдельно подчёркивает эту легко пропускаемую тонкость.

% ---------------------------------------------------------------------------
\section*{Итоги}
\begin{itemize}
  \item В Tor DoS-атаки одновременно являются атаками на анонимность: принудительное отключение или отказ реле может перенаправлять трафик к противнику.
  \item Память защищает внутренний лимит \opt{MaxMemInQueues}; при его превышении tor очищает кэши и закрывает цепи, накапливающие данные.
  \item Входные узлы применяют к клиентам ограничения соединений и цепей по IP-адресам с помощью корзин токенов. Отключённые по умолчанию в коде механизмы включаются и регулируются консенсусом.
  \item Onion-сервисы сталкиваются с амплификацией: инициировать запрос дёшево, обслуживать его дорого.
  \item Доказательство выполнения работы v1 сочетает Equi-X с проверкой Blake2b, обеспечивающей линейный рост усилия. Сервис отдаёт приоритет запросам с большим усилием и регулирует рекомендуемое усилие по состоянию очереди.
\end{itemize}

% ---------------------------------------------------------------------------
\section*{Упражнения}

\begin{esercizio}[1]\label{es:15-bucket}
При \texttt{DoSCircuitCreationRate} 2 и \texttt{Burst} 60, как в публичном консенсусе, сколько цепей может создать клиент за одну минуту, начиная с полного запаса токенов? А за один час? Считайте, что защита уже активирована, то есть адрес достиг \texttt{DoSCircuitCreationMinConnections} одновременных соединений. Ниже этого порога ограничения нет.
\end{esercizio}

\begin{esercizio}[1]\label{es:15-guard}
Почему ограничения по IP-адресам применяются только к клиентам и только на первом переходе? Что произойдёт, если промежуточный узел попытается ограничивать цепи \enquote{по адресу}?
\end{esercizio}

\begin{esercizio}[2]\label{es:15-amplificazione}
Оцените отношение работы клиента к работе сети для одного запроса введения: сколько цепей, рукопожатий и ячеек порождает ответ сервиса?
\end{esercizio}

\begin{esercizio}[2]\label{es:15-bit}
Сравните доказательство выполнения работы в Tor со схемой, требующей $k$ старших нулевых битов хеша. Какие уровни усилия допускает каждая схема? Почему линейная шкала лучше подходит для управления с обратной связью?
\end{esercizio}

\begin{esercizio}[2]\label{es:15-sforzo}
Клиент объявляет усилие $E = 50$, но случайно сразу находит решение с очень маленьким $R$, которое удовлетворяло бы и $E = 5000$. Может ли он отправить это решение, заявив усилие $E = 5000$? Почему?
\end{esercizio}

\begin{esercizio}[3]\label{es:15-oom}
В лабораторной сети установите \opt{MaxMemInQueues} в минимально допустимое значение на промежуточном реле и постройте через него цепь, загружая большой файл \emph{без чтения} входящих данных и без отправки SENDME. Что происходит с накопившимися ячейками? Найдите сообщения системы управления памятью в журнале реле.
\end{esercizio}

\begin{esercizio}[3]\label{es:15-coda}
Напишите небольшой симулятор очереди приоритетов onion-сервиса. Запросы поступают со случайным усилием, многие с нулевым и некоторые с более высокими значениями; сервис обрабатывает определённое число запросов в секунду и при превышении порога очереди отбрасывает половину с наименьшим усилием. Измерьте среднее время ожидания в зависимости от усилия.
\end{esercizio}

% ---------------------------------------------------------------------------
\section*{Дополнительная литература}
\begin{itemize}
  \item \titolo{Tor denial-of-service mitigations} (dos-spec) и \titolo{Onion service proof-of-work} (hspow-spec)~\cite{torspec}.
  \item R.~Jansen, F.~Tschorsch, A.~Johnson, B.~Scheuermann, \titolo{The Sniper Attack: Anonymously Deanonymizing and Disabling the Tor Network}~\cite{jansen2014sniper}.
  \item Репозиторий Equi-X tevador и журнал разработки, на которые ссылается спецификация.
\end{itemize}
